\documentclass[12pt]{article}
\pdfoutput=1

\usepackage{draft} 
\usepackage{hyperref}
\usepackage{graphicx,color,subfig}
\usepackage{cite}
\usepackage{mciteplus}
\usepackage{skak}
\usepackage{empheq}
\usepackage{anyfontsize}
% Use Chancery Font
\DeclareFontFamily{OT1}{pzc}{}
\DeclareFontShape{OT1}{pzc}{m}{it}{<-> s * [1.10] pzcmi7t}{}
\DeclareMathAlphabet{\mathpzc}{OT1}{pzc}{m}{it}

%\usepackage{chngcntr}
%\counterwithout{equation}{section}

\newcommand{\wt}[1]{\widetilde{#1}}
\newcommand{\vv}[1]{\left\langle #1 \right\rangle}
\newcommand{\un}[1]{\underline{#1}}
\newcommand{\BL}[1]{{ {\color{blue}[#1]}}}


\newcommand{\fixme}[1]{{\bf {\color{red}[#1]}}}
\def\be#1\ee{\begin{align}#1\end{align}}
\newcommand\nn{\nonumber}

\begin{document}

 
\unitlength = .8mm


\section{3-pt vertex}
Consider the light cone 3-pt vertex between three tachyon states. Let the incoming momenta (in $d-2$ transverse directions) be $p, p_1, p_2$ and the string lengths $p^+ = p_1^+ +p_2^+$. The corresponding tachyon wavefunctional defined on the unit disk takes the form
\ie
\Psi[p^i; X^{i}] \propto \exp \left[ -\frac{1}{\A'} \sum_m m X^i_m  X^i_{-m} + i p^i  X^i_0 \right], \qquad i = 1, \hdots, d-2.
\fe
Discretizing three unit circles into $L$ positions $z_k = e^{i 2 k \pi / L}$ for the first, and $L^{(r)} = \frac{p_r^+}{p^+} L$ positions for the second and third unit circles, we can then consider the product of the wavefunctionals, with the identification imposed \fixme{for now we assume that the singular point where the boundaries of all three cylinders touch is not one of the discretization points}:
\ie
&X^i(k) = X^{(1),i}(k), \text{ for } k = 1, \hdots, L^{(1)}\\
&X^i(k) = X^{(2),i}(k-L_1), \text{ for } k = L^{(1)} + 1, \hdots, L^{(1)} + L^{(2)} = L,
\fe
with the Fourier modes
\ie
X^i_m = \frac{1}{L} \sum_{k=1}^{L} X^i(k) e^{2 \pi i \frac{ k m}{L} },
\fe
and dispersion relation
\ie
m \to \frac{L}{\pi} \sin( \frac{\pi m}{L} ),
\fe
appropriately defined in the discretization.\\
\\
The 3-pt vertex is then given by
\ie
\propto \int& \prod_{i=1}^{d-2} \left(\left[\prod_{k=1}^{L} dX^i(k) \right]\prod_{r=1,2}\left[\prod_{m= 1}^{L^{(r)} } dX^{(r),i}(k) \right]\right) \delta_{\text{identification} }\\
\exp\Bigg[ & - \frac{1}{\pi \A' L }  \bigg( \sum_{m=1}^{L} \sin \frac{\pi m}{L}  \sum_{k,\ell}^L e^{2 \pi i \frac{(k-\ell) m}{L} } X(k) \cdot X(\ell) \bigg) \\
& - \frac{1}{\pi \A'} \sum_{r=1,2}  \bigg( \sum_{m=1}^{L^{(r)} }\frac{1}{L^{(r)}} \sin \frac{\pi m}{L^{(r)} }  \sum_{k,\ell}^{L^{(r)} } e^{2 \pi i \frac{(k-\ell) m}{L^{(r)} } } X^{(r)}(k) \cdot X^{(r)}(\ell) \bigg)\\
+&i \bigg( \frac{p}{L} \cdot \sum_k^L X(k) + \sum_{r=1,2} \frac{p_r}{L^{(r)}} \cdot \sum_k^{L^{(r)} } X^{(r)}(k) \bigg)
\Bigg].
\fe
After applying the identification, we can consider how the zero mode $X_0$ appears, which will only be in the linear term, which we can rewrite as
\ie
i \frac{(p + p_1 + p_2)}{L} \cdot  \sum_k^L X(k) + i \frac{(L^{(2)} p_1 - L^{(1)} p_2)}{L} \cdot\left(   \frac{1}{L^{(1)} }\sum_{k=1}^{L^{(1)}} X(k) - \frac{1}{L^{(2)} }\sum_{k=L^{(1)}+ 1}^{L} X(k)\right),
\fe
the first term being precisely proportional to the zero mode $X^i_0 = \frac{1}{L} \sum_{k=1}^L X(k) $ and generating the momentum conservation for transverse momenta. The second term is orthogonal to the zero mode, since 
\ie
\frac{1}{L^{(1)} }\sum_{k=1}^{L^{(1)}} 1 - \frac{1}{L^{(2)} }\sum_{k=L^{(1)}+ 1}^{L} 1 = 0,
\fe
and will contribute to the Gaussian integration. We can write the 3-pt vertex as
\ie
\propto \int \prod_{i=1}^{d-2} \left[ \prod_{k=1}^L dX^i(k) \right] \exp \left[ - \frac{1}{2 \pi \A'}  (X^i)^T \cdot M  \cdot X^i +i J^i \cdot X^i \right]e^{i (p + p_1 + p_2) \cdot X_0 }
\fe
with
\ie
%%% OLD FORMULA FOR SUM UP TO L/2, also before factoring out 1/\pi
%M_{k\ell} = &\frac{1}{\pi L}(\frac{-\sin \frac{\pi}{L} + (-1)^{k-\ell} (\cos \frac{\pi}{L} - e^{2 i \pi \frac{k- \ell}{L} })}{ \cos \frac{\pi}{L} - \cos \frac{2\pi (k-\ell)}{L} } ) \\
%+ \sum_{r=1,2}& \chi_r(k) \chi_r(\ell) \frac{1}{\pi L^{(r)} }(\frac{-\sin \frac{\pi}{L^{(r)}} + (-1)^{k_r-\ell_r} (\cos \frac{\pi}{L^{(r)}} - e^{2 i \pi \frac{k_r- \ell_r}{L^{(r)}} })}{ \cos \frac{\pi}{L^{(r)}} - \cos \frac{2\pi (k_r-\ell_r)}{L^{(r)}} } ) ,
M_{k\ell} = &\frac{1}{ L}(\frac{-\sin \frac{\pi}{L}}{ \cos \frac{\pi}{L} - \cos \frac{2\pi (k-\ell)}{L} } ) \\
+ \sum_{r=1,2}& \chi_r(k) \chi_r(\ell) \frac{1}{ L^{(r)} }(\frac{-\sin \frac{\pi}{L^{(r)}} }{ \cos \frac{\pi}{L^{(r)}} - \cos \frac{2\pi (k_r-\ell_r)}{L^{(r)}} } ) ,
\fe
where 
\ie
&k_1 \equiv k,\\
&k_2 \equiv k - L^{(1)}
\fe
and
\ie
\chi_r(k_r) = 
\begin{cases}
1, 1 \leq k_r \leq L^{(r)}\\
0, \text{ otherwise } 
\end{cases}
\fe
just tests if $k$ is in the appropriate range so that $X^i(k)$ can be identified with $X^{(r),i}(k)$. $M$ should have one zero eignevector corresponding precisely to the zero mode.  
And,
\ie
J_k^i = \frac{(L^{(2)} p_1^i - L^{(1)} p_2^i)}{L}  \left( \frac{\chi_1(k)}{L^{(1)} } - \frac{\chi_2(k)}{L^{(2)} } \right)
\fe
satisfies $\sum_k J_k^i = 0$.\\
\\
The 3-pt function then becomes
\ie
\propto ~~ (2\pi)^{d-2} \delta^{(d-2)}(p + p_1 + p_2)\, \left( {\det}' M   \right)^{-\frac{d-2}{2} } \,e^{-\frac{\pi \A'}{2} J^{i,T} \cdot {M}^{+} \cdot J^i },
\fe
where ${\det}'$ denotes product of all nonzero eigenvalues, and $M^+$ is the Moore-Penrose pseudoinverse
\ie
M^+ = \sum_{s=1}^{L-1} \frac{1}{\lambda_s} \xi_s \xi_s^T 
\fe
for $\xi_s$ the eigenvectors of $M$ with eigenvalues $\lambda_s \neq 0$. At this point it's also useful to be careful about the normalization. Considering the sphere 2-pt function in terms of the disk wavefunctionals, we should demand that 
\ie
\int \prod_{i=1}^{d-2}\left[ \prod_{k=1}^L dX^i(k) \right] \Psi[p^i; X^i] \Psi^{\text{c}}[q^i; X^i] = (2|p^+|) (2\pi)^{d-2} \delta^{(d-2)}(p^i+q^i)
\fe
with $\Psi^{\text{c}}$ the conjugate wavefunctional, where $X^i_n \leftrightarrow X^i_{-n}$. We note that the nonzero modes just give a Gaussian integral of $\exp(- \frac{1}{2\pi \A'} (X^i)^T \cdot M' \cdot X^i )$ where 
\ie
M'_{k \ell} =  \frac{2}{ L}(\frac{-\sin \frac{\pi}{L}}{ \cos \frac{\pi}{L} - \cos \frac{2\pi (k-\ell)}{L} } )
\fe
whose ${\det}' M' = \prod_{m=1}^{L-1}2  \sin \frac{\pi m}{L} = L$. We make the change of variables
%  $X^i(L) = LX_0 - \sum_{k=1}^{L-1} X(k)$ 
\ie
X(k) &= X_0 + \tilde{X}(k), \qquad k = 1,\hdots, L-1\\
 X(L) &= X_0 - \sum_{k=1}^{L-1} \tilde{X}(k)
\fe
and define the corresponding $(L-1) \times (L-1)$ matrix $\wt{M}'$ whose $\det \wt{M}' = L^2$. But the change of variables comes with its own Jacobian factor $L$ which precisely cancels $(\det \wt{M}')^{-\frac{1}{2}}$. So the normalization factor for $\Psi[p^i;X^i]$ on the disk with boundary discretized into $L$ steps is $\sqrt{2 |p^+|\,}\left((2\pi^2 \A')^{L-1} \right)^{-\frac{d-2}{4}}$. With the same change of variables we can define a reduced matrix and vector $\wt{M}_{k\ell}$ and $\wt{J}^i_k$ where $k$ now runs from 1 to $L-1$. Using the normalizations the 3-pt function becomes 
\ie
=\sqrt{8 |p^+ p_1^+ p_2^+|\,} (2\pi)^{d-2} \delta^{(d-2)}(p + p_1 + p_2)\, (2\pi^2 \A')^{ \frac{d-2}{4} }\left(\frac{ {\det} \wt{M}  }{L^2}    \right)^{-\frac{d-2}{2} } \,e^{-\frac{\pi \A'}{2} \wt{J}^{i,T} \cdot \wt{M}^{-1} \cdot \wt{J}^i }.
\fe
We note that $\det \wt{M} = L \,  {\det}' M$. The ratio $\frac{\det \wt{M}(L^{(1)}, L^{(2)}) }{\det \wt{M}(L^{(1)}{}', L^{(2)}{}')}$ converges quickly, but depends on $(L^{(1)}, L^{(2)})$. Numerically we checked that $\det \wt{M}$ itself seems to scale as $L^{9/4}$, i.e. ${\det}' M$ scales as $L^{5/4}$.\\
\\
Meanwhile, 
\ie
{\wt J}^i_k &= \frac{(L^{(2)} p_1^i - L^{(1)} p_2^i)}{L}  \left( \frac{\chi_1(k)}{L^{(1)} } - \frac{\chi_2(k)}{L^{(2)} } + \frac{1}{L^{(2)} }   \right)\\
&= \frac{(L^{(2)} p_1^i - L^{(1)} p_2^i)}{ L^{(2)} L^{(1)} } \{1,\hdots,1,0,\hdots,0\}\\
&= \frac{p^+}{L} \left(\frac{p_1^i}{p_1^+} - \frac{p_2^i}{p_2^+} \right) \{1,\hdots,1,0,\hdots,0\}
\fe
with $L^{(1)}$ ones and $(L^{(2)}-1)$ zeros. $\wt{J}^{i,T} \cdot \wt{M}^{-1} \cdot \wt{J}^i$ is proportional to the sum of all the terms in the first $L^{(1)} \times L^{(1)}$ block of $\wt{M}^{-1}$. We checked that $\wt{J}^{i,T} \cdot \wt{M}^{-1} \cdot \wt{J}^i = J^{i,T} \cdot {M}^{+} \cdot J^i $, and that its value
\ie
\wt{J}^{i,T} \cdot \wt{M}^{-1} \cdot \wt{J}^i=(p^+)^2  \left(\frac{p_1^i}{p_1^+} - \frac{p_2^i}{p_2^+} \right)^2 f( p_1^+, p_2^+,L)
\fe
with $f$ converging fairly rapidly, but depending explicitly on $p_1^+,\,p_2^+$. Using transverse momentum conservation, we can rewrite this as 
\ie
\wt{J}^{i,T} \cdot \wt{M}^{-1} \cdot \wt{J}^i = \frac{(p^+)^3}{p_1^+ p_2^+} \left( \frac{(p_1^i)^2}{p_1^+} +  \frac{(p_2^i)^2}{p_2^+} -  \frac{(p^i)^2}{p^+} \right) f(p_1^+, p_2^+,L),
\fe
and for on-shell tachyons $-\frac{4}{\A'} = m^2 = 2 p^+ p^- - (p^i)^2$ and this becomes
\ie
\label{eq:JMJ}
\wt{J}^{i,T} \cdot \wt{M}^{-1} \cdot \wt{J}^i =  \frac{(p^+)^3}{p_1^+ p_2^+}   f( p_1^+, p_2^+,L) \left(2(p_1^- +p_2^- - p^-) + \frac{4}{\A'} (\frac{1}{p_1^+} + \frac{1}{p_2^+} - \frac{1}{p^+}) \right).
\fe
Putting everything together, the 3-pt vertex becomes (assuming on-shell momenta and using conservation of $p^+$ and $p^i$:
\ie
\label{eq:3ptvertex}
= ~~ & \sqrt{8 |p^+ p_1^+ p_2^+|\,} (2\pi^2 \A')^{\frac{d-2}{4} }\delta_{p^+,  p_1^+ + p_2^+} (2\pi)^{d-2} \delta^{(d-2)}(p^i + p_1^i + p_2^i) e^{- \pi \A' \frac{(p^+)^3}{p_1^+ p_2^+} f(p_1^+, p_2^+, L) (p_1^- +p_2^- - p^-)} \\
& ~~ \left(\frac{\det \wt{M}(p_1^+,p_2^+,L) }{L^2 } \right)^{-\frac{d-2}{2} }  \exp\left(-\pi \frac{(p^+)^2}{(p_1^+ p_2^+)^2} f(p_1^+, p_2^+, L) ((p^+)^2 + (p_1^+)^2 + (p_2^+)^2)  \right)
\fe
Note that $\delta_{p^+,  p_1^+ + p_2^+} $ is imposed at the level of the discretization scheme. A key test is that we should recover full Lorentz invariance \fixme{of precisely what expression?} for $d=26$ as $L \to \infty$. Conservation of energy, once it is recovered from integrating over the time $x^+$ at which the interaction occurs, should set the first exponential to 1.
%Therefore we might expect for $d=26$ that the expressions on the second line evaluate to a constant (for fixed $L$ and varying $p_1^+, \, p_2^+$) \fixme{but numerically this doesn't seem to work yet}.
\\
\\
We note that one should consider the contribution of the Weyl anomaly at the interaction point. From (2.4) in [1102.2751] the Weyl anomaly factor is (let $x \equiv \frac{p_1^+}{p^+}$ and $1-x = \frac{p_2^+}{p^+} $):
\ie
e^{-\frac{d-2}{24}\Gamma[\log (\partial \rho \bar \partial \rho)]} \equiv \left(\frac{\exp [-2 \hat{\tau}_0 \sum_r \frac{1}{\alpha_r} ] }{ \prod_r \alpha_r  } \right)^{\frac{d-2}{24} } = \left(\frac{\exp [-2 \frac{\hat{\tau}_0}{2p^+} (1 - \frac{1}{x} - \frac{1}{1-x} ) ] }{(2 p^+)^3 x(1-x)  } \right)^{\frac{d-2}{24} }
\fe
where $\frac{d-2}{24}$ \fixme{came from chatGPT, in the paper they are using d=26 so it is just 1},  
\ie
\hat{\tau}_0 \equiv \sum_r \alpha_r \log|\alpha_r| = 2p^+ (-x \log x - (1-x) \log(1-x) ),
\fe
$\rho(z)$ is the Mandelstam mapping of the Mandelstam diagram to the complex plane, and in our conventions $\alpha_3 = 2p^+$, $\alpha_1 = -2p_1^+=-2p^+ x$, $\alpha_2 = -2p_2^+= -2p^+(1-x)$.\\
We find that, as functions of $x$,
\ie
{}&-\Gamma[\log (\partial \rho \bar\partial \rho)](x) - 3\log(2p^+)\\
 &\qquad= \log \left[ \frac{\exp [2 (x \log x + (1-x) \log(1-x)) (1 - \frac{1}{x} - \frac{1}{1-x} ) ] }{ x(1-x)  } \right] \\
&\qquad=\text{const} - 12 \lim_{L\to\infty}\log \frac{\det \wt{M}(x,L)}{L^{9/4} }
\fe
and
\ie
-\frac{1}{\pi} x(1-x) \left( x \log x + (1-x) \log(1-x) \right) = \lim_{L\to\infty} f(x,L).
\fe
Comparing to the last exponential in (\ref{eq:3ptvertex}), or the second term in (\ref{eq:JMJ}), we see that it takes the form
\ie
-2\pi &\frac{1}{x(1-x)} f(x,L) \left(\frac{1}{x}+\frac{1}{1-x}-1 \right) \\
&\qquad\to -2\left( x \log x + (1-x) \log(1-x) \right)\left(1-\frac{1}{x}-\frac{1}{1-x} \right),
\fe
which precisely cancels the exponential piece of $e^{-\Gamma}$ coming from $(\frac{\det \wt{M}}{L^{9/4} })^{-\frac{d-2}{2} } =C e^{- \frac{d-2}{24} \Gamma }$, provided $d=26$!\\
\\
inputting the conjectured forms of $f(x,L)$ and $\det \wt{M}(x,L)$ for $L \to \infty$, we get the 3-pt vertex
\ie
\label{eq:3ptvertexConj}
=&\sqrt{8 |p^+|^3\,} C\,\frac{(2\pi^2 \A')^6}{L^3} \frac{1}{ \sqrt{x(1-x) } } \delta_{p^+, p_1^+ +  p_2^+}  (2\pi)^{24} \delta^{(24)}\Big(p^i + p_1^i + p_2^i\Big)\\
&~~ \exp\left[\A' p^+(p_1^- +p_2^- -p^-) (x \log x + (1-x) \log (1-x)) \right].
\fe
Note that when calculating the \textit{amplitude}, one must integrate over LC time $x^+$, giving a $\delta(p^- -  p_1^- - p_2^-)$ that kills the exponential term.  \fixme{Check that this works out, I still hacen't hunted down all the normalizations} In addition, the $\frac{1}{\sqrt{x(1-x)} }$ precisely corresponds to $\prod_r (\A_r)^{-1/2}$ in GSW (11.2.7). The remaining issue is to explain the factor $\frac{1}{L^3}$. In fact $ L \delta_{p^+,p_1^+ + p_2^+} \to \delta(p^+ - p_1^+ - p_2^+) $, so it is an overall factor of $\frac{1}{L^4}$ that must be explained.
%
%
%\fixme{The first factor above, including the $\frac{1}{L}$'s should somehow be related to integrating over $dp^+ dp_1^+ dp_2^+$ with the appropriate measure factor. For example in [[1102.2751] (2.2), each $dp_r^+$ comes with a factor of $p_r^+$, which would cancel against the $x(1-x)$ in the denominator, which is the last remaining $x$ dependence. The $L$'s should turn the sums into integrals because for a fixed $p^+$ and $L$, $x$ can only take discrete values $x = \frac{1}{L}, \frac{2}{L}, \hdots \frac{L-1}{L}$ with $\Delta p_1^+ = \frac{p^+}{L} \to dp_1^+$ as $L \to \infty$. But then it seems there's too many $L$'s in the denominator.  }

\section*{Normalization conventions and comparison with GSW}
For on-shell external strings in light-cone variables, the Lorentz-invariant one-particle measure is
\ie
d\Pi(p)=\frac{d^{26}p}{(2\pi)^{25}}\delta(p^2+m^2)\theta(p^0)=\frac{dp^+\,d^{24}p^i}{(2\pi)^{25}\,2p^+}.
\fe
If we first use delta-function normalized light-cone momentum eigenstates,
\ie
{}_0\langle p'|p\rangle_0=(2\pi)^{25}\delta(p'^+-p^+)\delta^{(24)}(p'^i-p^i),
\fe
then the relativistically normalized states are
\ie
|p\rangle_{\rm rel}=\sqrt{2p^+}\,|p\rangle_0,
\fe
so that
\ie
{}_{\rm rel}\langle p'|p\rangle_{\rm rel}=2p^+(2\pi)^{25}\delta(p'^+-p^+)\delta^{(24)}(p'^i-p^i).
\fe
This is precisely the normalization imposed above in the two-point function, so each external disk wavefunctional already carries a factor $\sqrt{2|p_r^+|}$.

Therefore if $A_3$ denotes the amputated Lorentz-invariant three-point amplitude in the convention of GSW, while $S_3$ denotes the three-point matrix element between relativistically normalized external states, then
\ie
S_3=(2\pi)^{26}\delta^{(26)}(p_1+p_2-p_3)\,\frac{A_3}{\sqrt{(2p_1^+)(2p_2^+)(2p_3^+)}} ,
\fe
equivalently
\ie
A_3=\sqrt{(2p_1^+)(2p_2^+)(2p_3^+)}\,\frac{S_3}{(2\pi)^{26}\delta^{(26)}(p_1+p_2-p_3)}.
\fe
In our conventions
\ie
\alpha_3=2p^+,\qquad \alpha_1=-2p_1^+,\qquad \alpha_2=-2p_2^+,
\fe
so the GSW external momentum dressing is
\ie
\prod_{r=1}^3 |\alpha_r|^{-1/2}=\frac{1}{\sqrt{|\alpha_1\alpha_2\alpha_3|}}=\frac{1}{\sqrt{8|p^+p_1^+p_2^+|}}=\frac{1}{(2p^+)^{3/2}\sqrt{x(1-x)}}.
\fe
Thus the factor $\frac{1}{\sqrt{x(1-x)}}$ appearing in (\ref{eq:3ptvertexConj}) reproduces the $x$-dependence of the GSW leg factor, but the full GSW factor also contains an overall $(2p^+)^{-3/2}$. Since our wavefunctionals were normalized to relativistic external states from the start, the quantity computed here should be compared directly with $S_3$, not with the amputated $A_3$. The remaining mismatch in the overall powers of $p^+$ and in the factors of $L$ should therefore come from the longitudinal normalization and the conversion of the discrete $p^+$ sums and Kronecker deltas to continuum integrals and Dirac delta functions, rather than from any additional external momentum dressing.
\end{document}

